\subsection{Historical context and the role of the method}

Vertex Cover appeared as \textsc{Node Cover} among the problems
Karp proved NP-complete in 1972~\cite{Karp1972}. Exact hardness
still leaves room for approximation, and the elementary
maximal-matching algorithm achieves factor two. Subsequent algorithms
improved this ratio by an amount depending on the graph order.
Karakostas obtained \(2-\Theta(1/\sqrt{\log n})\) using semidefinite
programming~\cite{Karakostas2009}. This saving tends to zero as
\(n\) grows. The threshold addressed here concerns a fixed saving
below two, so it is compatible with such algorithmic improvements.

The connection between proof checking and approximation hardness
was developed by Feige, Goldwasser, Lov\'asz, Safra, and
Szegedy~\cite{FGLSS96}. Their graph encodes accepting local proof
views, with consistency of the answers governing adjacency.
The PCP breakthrough of 1992, developed by Arora and Safra and
completed in constant-query form by Arora, Lund, Motwani, Sudan,
and Szegedy~\cite{AroraSafra98,ALMSS98}, yielded constant
approximation gaps for problems including Vertex Cover. Raz's
parallel repetition theorem~\cite{Raz1998} reduces the soundness
of projection games to any desired fixed positive constant while
preserving perfect completeness. Together these results supply the
Label Cover promise stated precisely in \Cref{src:hardness}, the
only external hardness input used here.

For Vertex Cover, H{\aa}stad established hardness for every fixed
factor below \(7/6\)~\cite[Theorem~8.1]{Hastad2001}.
Dinur and Safra raised this to every fixed factor below
\(10\sqrt5-21\), approximately \(1.36068\)~\cite[Theorem~1.1]{DinurSafra2005}.
Their analysis uses biased set families and Boolean-function
structure to control independent-set density. Khot's Unique Games
Conjecture~\cite{Khot02} proposed a stronger source of approximation
gaps. Assuming this conjecture, Khot and Regev obtained hardness
for every fixed factor below two~\cite{KhotRegev2008}. Their proof
uses a biased long code and an equivalent, stronger formulation
of the source promise that supports short lists of candidate labels.
Friedgut's approximation theorem for Boolean functions of low
average sensitivity~\cite{Friedgut1998} supplies bounded sets of
relevant coordinates in that analysis.

The 2017--2018 developments in the \(2\)-to-\(2\) Games program
gave an unconditional advance. The Grassmann-graph reduction of Khot, Minzer, and
Safra~\cite{KhotMinzerSafra2025}, the work of Dinur, Khot, Kindler,
Minzer, and Safra~\cite{DKKMS2025}, and the contribution of Barak,
Kothari, and Steurer~\cite{BarakKothariSteurer2019} culminated in
the expansion theorem of Khot, Minzer, and
Safra~\cite{KhotMinzerSafra2023}. This completed the Games Theorem
with imperfect completeness and yielded unconditional Vertex Cover
hardness for every fixed factor below \(\sqrt2\). The resulting
independent-set gap has completeness density close to
\(1-1/\sqrt2\) and arbitrarily small soundness density. The gap
in \Cref{thm:gap} instead has completeness density approaching
\(1/2\), which is what gives the factor-two threshold.

Our argument retains the objective of extracting short local label
lists from a large independent set, but obtains the lists from
Lipschitz derivatives. The central step preserves variance while
restricting the information available to each list. We adapt the
product-space orthogonal decomposition of the Efron--Stein
method~\cite[Section~2]{EfronStein1981}, combining it with a
resampling estimate that controls the loss caused by this information
restriction. A variance inequality on a product of intervals then
supplies gradient energy, from which short candidate-label lists can
be extracted without an alphabet-size factor. These ingredients
are proved in full in \Cref{sec:restriction,ana:private-section};
the final decoding uses the ordinary Label Cover promise above.
